\documentclass[11pt]{article}
\usepackage[margin=25mm]{geometry}
\usepackage{fontspec}
\setmainfont{Latin Modern Roman}
\usepackage{amsmath,amssymb,mathtools}
\usepackage{unicode-math}
\setmathfont{Latin Modern Math}
\usepackage{xurl,hyperref}
\hypersetup{colorlinks=true,urlcolor=blue}
\setcounter{secnumdepth}{0}
\setlength{\parindent}{0pt}
\setlength{\parskip}{5pt}
\setlength{\emergencystretch}{3em}
\newcommand{\sref}[2]{\hyperlink{source:#1:#2}{[S#2]}}
\newcommand{\cataloguescope}[1]{\paragraph{Recorded target.}#1}
\newcommand{\R}{\mathbb R}
\newcommand{\C}{\mathbb C}
\newcommand{\N}{\mathbb N}
\newcommand{\Z}{\mathbb Z}
\newcommand{\Q}{\mathbb Q}
\newcommand{\F}{\mathbb F}
\newcommand{\eps}{\varepsilon}
\begin{document}
% UnsolvedMath ID 30004880; source catalogue code OWR-8415354-006
\setcounter{section}{30004879}
\section{Trigonometric and Elliptic R-Matrices from Three-Dimensional WZW Models}\label{30004880}
{\small UnsolvedMath ID 30004880 · Mathematical Physics}
\subsection{Definitions and mathematical statement}

\paragraph{Shared notation.}$\mathbb N=\{1,2,\ldots\}$; the standard number fields have their usual meanings. All additional conventions are stated locally.


An $R$-matrix acts on a tensor product and obeys a Yang--Baxter compatibility equation. Rational, trigonometric and elliptic refer to its dependence on a spectral parameter. In the report's three-dimensional Wess--Zumino--Witten construction, the chosen differential form is part of the action and controls the resulting integrable structure.

\paragraph{Catalogue wording retained for traceability.}
Can three-dimensional Wess–Zumino–Witten models be generalized to produce trigonometric and elliptic R-matrices for broader choices of the defining differential form?

\paragraph{Scope and source qualification.}The source asks for an extension of its particular construction. It does not assert that every arbitrary differential form produces an integrable theory.
\subsection{Short English statement}
Extend the three-dimensional WZW construction to broader differential-form choices producing trigonometric or elliptic R-matrices.
\subsection{Research attempt}
\paragraph{Model, date and investigated scope.}
OpenAI GPT-6 Astra (configured model \texttt{gpt-6-astra}), ultra reasoning; generated 11 October 2026. Derived a trigonometric local-pole kernel and an elliptic single-pole obstruction; no WZW construction or Yang–Baxter R-matrix.

\paragraph{Sources and literature check.}
Read the supplied definition, catalogue wording, source qualifications and reference entries for 30004880. The historical source/context and inherited status are retained; no fresh literature or novelty certification is claimed. No external theorem is used without proof in this attempt; explicit model assumptions are distinguished from facts about the source model.

\paragraph{A spectral-geometry test.} Passing from an additive coordinate \(u\) to \(z=e^u\) changes the logarithmic differential by \(dz/z=du\). The simple meromorphic kernel
\[
 K(u)=\frac1{e^u-1}+\frac12=\frac12\coth(u/2)
\]
is \(2\pi i\)-periodic and has expansion \(K(u)=u^{-1}+O(u)\) near zero. This supplies a concrete way rational local pole behavior can coexist with trigonometric global dependence after changing the underlying spectral geometry.
\paragraph{An elliptic obstruction.} A doubly periodic meromorphic scalar function cannot have just one simple pole with nonzero residue in a period cell. Integrate it around a translated fundamental parallelogram avoiding poles. Opposite-edge integrals cancel by periodicity, whereas a small contour around the sole pole gives \(2\pi i\) times its residue. The residue must therefore vanish, a contradiction. A naive elliptic replacement of the one-pole scalar kernel must introduce compensating poles, nontrivial quasi-periodicity or bundle-valued structure.
\paragraph{Source gap.} Neither kernel calculation constructs the source three-dimensional WZW action or a quantum \(R\)-matrix. The Yang--Baxter equation, gauge consistency, boundary terms and form-dependent quantization remain to be derived. The first gap is showing that a permissible new differential form and global spectral geometry produce the required operator exchange relations, including elliptic compensations. Arbitrary forms are not assumed to give integrability. These elementary analytic checks only delimit a possible extension route.

\paragraph{Verification and adversarial audit.}
Checked the logarithmic coordinate, local pole normalization and opposite-edge residue cancellation; distinguished scalar kernels from matrix exchange solutions. This is a GPT-6 Astra Ultra self-audit, not independent human verification. Prompt A drafting followed by prompt B scope/dependency/boundary self-audit. Checked the logarithmic coordinate, local pole normalization and opposite-edge residue cancellation; distinguished scalar kernels from matrix exchange solutions. The remaining gap is: Derive consistent quantum exchange relations from the source WZW action for new global form/geometry choices. Unresolved status and zero original-target completion are retained.

Coordinator checked the logarithmic coordinate kernel and the residue obstruction for a doubly periodic one-pole scalar function. This additional check was model-assisted, not independent human verification.

\paragraph{Tools, process and contributions.}
Python 3 and local file-reading/writing tools for manuscript serialization and structural checks; mathematical calculations were symbolic reasoning, with no numerical experiment or theorem prover. Contribution prompt A was followed by prompt B. GPT-6 Astra with ultra reasoning drafted the research and performed its stated self-audit; any additional coordinator contribution is identified in the verification and audit. A GPT-6-based Codex coordinator assembled the manuscript and checked repository integration; its exact variant and reasoning configuration were not exposed. No human mathematical contribution, independent human proof check or formal proof certification is supplied. Complete derivations and the consulted source references are included here; private reasoning transcripts are not published.

\paragraph{Final status and first remaining gap.}
\textbf{UNRESOLVED}. Derive consistent quantum exchange relations from the source WZW action for new global form/geometry choices.

\section{Completion Estimate}
\noindent\textbf{COMPLETION: 0\%}
The proved calculation is an elementary auxiliary result or a test of a restricted strategy; no new portion of the full source target is established.

\subsection{Sources}
\begin{itemize}
\item[\hypertarget{source:30004880:1}{S1}] ULAM.AI and the UnsolvedMath Contributors, UnsolvedMath, record 30004880 (OWR-8415354-006). Catalogue curation under CC BY 4.0; definitions and English exposition added locally.
\url{https://huggingface.co/datasets/ulamai/UnsolvedMath}.
\item[\hypertarget{source:30004880:2}{S2}] Mini-Workshop: Three Facets of R-Matrices (2021), Oberwolfach Reports, DOI 10.4171/owr/2021/51, pp. 2791--2825. Original report and mathematical context.
\url{https://doi.org/10.4171/owr/2021/51}.
\end{itemize}
\label{end:30004880}
\end{document}
